\documentclass[11pt,a4paper]{article}
\usepackage[utf8]{inputenc}
\usepackage[T1]{fontenc}
\usepackage{amsmath,amssymb,amsthm,geometry,hyperref,booktabs}
\geometry{margin=1in}
\newtheorem{theorem}{Theorem}
\newtheorem{lemma}{Lemma}
\newtheorem{proposition}{Proposition}
\newtheorem{conjecture}{Open problem}
\newcommand{\sphere}{\mathbb{S}^{2}_{R}}
\title{V5-32: Constructive Weak Pentagonal Refinement\\
on the Sphere and Topological Obstructions\\
{\large A Research Preprint Draft}}
\author{Long Ngo\\
Vietflexmap\\
\small Copyright (c) 2025 Vietflexmap \quad \small Long Ngo 2026}
\date{October 2026\\{\small Apache License, Version 2.0}}
\begin{document}
\maketitle
\begin{abstract}
We establish an elementary constructive existence theorem for weak nested
32-ary pentagonal partitions of the sphere. Starting with twelve planar
pentagonal faces of a convex dodecahedron, a two-piece nonconvex pentagonal
refinement can be iterated five times to produce 32 strict pentagonal
descendants per parent. Radial projection transfers the nested construction
to the sphere, yielding exactly $12\cdot32^r$ spherical pentagons at each
finite resolution $r$. The weak construction permits nonconvexity and
hanging nodes; it does not prove edge-to-edge conformity, uniform
shape-regularity, equal-area refinement, point-location efficiency, or
performance superiority to existing DGGS. We derive necessary combinatorial
constraints from Euler's relation, an edge-count parity obstruction, and a
restriction on fivefold rotational symmetry for convex 32-child refinements.
We give a deterministic exact-rational prototype and separate the established
claims from open problems for a strong V5 grid.
\end{abstract}
\noindent\textbf{Keywords:} pentagonal tilings, spherical partitions,
hierarchical grids, combinatorial topology, discrete global grid systems,
constructive geometry.\\
\noindent\textbf{Repository:} \url{https://github.com/Base27-CVNSS/V5}

\section{Introduction}
Geospatial indexing schemes such as H3 organize coordinates into hierarchical
cell identifiers. H3 is predominantly hexagonal and has twelve exceptional
pentagons; its implementation and indexing scheme are independent of V5
\cite{h3docs,h3code}. Here we study a different, narrowly defined existence
question: can a pentagonal parent be recursively divided into exactly 32
pentagonal children, and can a family of such partitions cover the sphere?
It is essential to distinguish a weak geometric partition from a production
discrete global grid system (DGGS), whose requirements include efficient
geospatial operations and usable geometric quality \cite{ogc}.

\section{Definitions}
A \emph{strict pentagon} is a simple planar polygon with exactly five distinct
vertices and five edges, with no consecutive triple collinear. Convexity is
not required. A \emph{weak partition} of a region is a finite family of cells
with disjoint interiors whose union is the region; boundary hanging nodes
are permitted. An \emph{edge-to-edge} partition instead requires all shared
boundaries to be complete common edges or vertices.

\section{Constructive existence}
\begin{lemma}[Bent-chord binary refinement]
Every strict simple pentagon admits a weak subdivision into two strict
simple pentagons.
\end{lemma}
\begin{proof}
Triangulate $P=(A,B,C,D,E)$. The two diagonals of a triangulation of a
pentagon have a common endpoint; relabel such that they are $AC$ and
$AD$, giving triangles $ABC$, $ACD$, and $ADE$. Choose
$Q\in\operatorname{relint}(CD)$, so the chord $AQ$ runs inside triangle
$ACD$. Choose $T$ sufficiently close to an interior point of $AQ$,
strictly inside $ACD$, avoiding the line $AQ$ and the finitely many
lines causing consecutive vertices to become collinear.
The proper crosscut $A$--$T$--$Q$ separates the domain into
$(A,B,C,Q,T)$ and $(A,T,Q,D,E)$. Both are simple pentagons
with strict corners. Their interiors are disjoint and their union is $P$.
One of them may be nonconvex.
\end{proof}

\begin{theorem}[Weak 32-ary refinement]
For every finite integer $r\geq0$, every strict simple pentagon admits
a weak nested partition into $32^r$ strict pentagons.
\end{theorem}
\begin{proof}
At each refinement step apply the lemma to every current pentagon.
Five binary steps produce $2^5=32$ children. Recursively repeating
this five-step operation $r$ times yields $32^r$ descendants while
preserving exact set coverage and disjoint interiors.
\end{proof}

\begin{theorem}[Global spherical existence]
For every finite integer $r\geq0$, the sphere $\sphere$ admits a
weak nested partition into $12\cdot32^r$ strict spherical
pentagons bounded by great-circle arcs.
\end{theorem}
\begin{proof}
Let $D$ be a regular convex dodecahedron with its center at the
origin. Its boundary has twelve pentagonal faces. Apply the preceding
refinement to each face, constructing $12\cdot32^r$ planar cells in
total. The radial map $\rho:\partial D\to\sphere$ given by
$\rho(x)=Rx/\lVert x\rVert$ is a homeomorphism: convexity and the
interior origin ensure every radial ray intersects $\partial D$
once, and compactness establishes a continuous inverse. It preserves
coverage and disjoint interiors. Within a face plane not passing
through the origin, every straight edge and the origin span a plane,
whose intersection with $\sphere$ is a great circle. The restriction
of $\rho$ to a face is a smooth local diffeomorphism and preserves
the noncollinearity of consecutive edge directions. Thus each cell
is a strict spherical pentagon with five geodesic arcs.
\end{proof}

\section{Necessary restrictions under edge-to-edge conformity}
\begin{proposition}[Euler degree budget]
In an edge-to-edge pentagonal decomposition of a sphere into $F$
faces, where $d(v)$ is the vertex degree,
\begin{equation}
\sum_v (d(v)-3) = \frac{F}{2}-6.
\end{equation}
For $F=12\cdot32^r$, the right-hand side equals $6(32^r-1)$.
\end{proposition}
\begin{proof}
Five incidences per face and two per edge imply $5F=2E$.
Euler's equation $V-E+F=2$ gives $V=3F/2+2$.
Finally $\sum_vd(v)=2E=5F$, from which the claim follows.
\end{proof}

\begin{proposition}[Boundary parity]
If an edge-to-edge partition of a pentagonal parent consists of
$F$ pentagonal children with $I$ interior edges and $B$ boundary
edges, then $B\equiv F\pmod2$. In particular $F=32$ requires an
even number of boundary segments; retaining exactly five unsplit
parent sides is impossible.
\end{proposition}
\begin{proof}
Count cell-edge incidences: $5F=2I+B$.
\end{proof}

\begin{proposition}[Fivefold symmetry obstruction]
A partition of a fivefold-rotationally symmetric convex pentagon
into $N$ convex polygonal cells invariant as a set under rotation
by $2\pi/5$ must satisfy $N\equiv0$ or $1\pmod5$.
Consequently $N=32$ is impossible under these combined assumptions.
\end{proposition}
\begin{proof}
The rotation group $C_5$ acts on cells. Each orbit has size $1$
or $5$. An invariant full-dimensional convex cell contains the
rotation center in its interior. At most one child can contain the
center in its interior, because interiors are pairwise disjoint.
Thus $N=5k$ or $N=5k+1$.
\end{proof}

\begin{proposition}[Regular spherical cell obstruction]
An edge-to-edge tiling of the sphere by congruent convex regular
spherical pentagons with geodesic edges has exactly twelve faces.
\end{proposition}
\begin{proof}
If the common interior angle is $\alpha$, Girard's formula gives
$A/R^2=5\alpha-3\pi>0$, implying $\alpha>3\pi/5$.
Convexity gives $\alpha<\pi$. The local angle sum at any vertex
requires $d\alpha=2\pi$ with integer $d\geq3$, hence $d=3$.
Substituting into the Euler degree budget yields $F=12$.
\end{proof}

\section{Computational prototype}
The standard-library Python prototype at \url{https://github.com/Base27-CVNSS/V5/blob/main/src/v5.py}
uses exact rational arithmetic for point construction, signed
areas and strict-corner checks. Unit tests assert exact area conservation
and the correct 32- and 1024-cell cardinalities. Such checks are
illustrative, rather than independent certification of global
intersection-freedom.

\section{Open problems and limitations}
\begin{conjecture}[Strong 32-ary V5, open]
Find a globally edge-to-edge, recursively 32-ary spherical
pentagonal grid with convex cells, a uniform shape-regularity
constant, and $\max_C \mathrm{diam}(C)=O(32^{-r/2})$.
\end{conjecture}
No proof is presented for this conjecture. The weak constructive
theorem does not establish shrinking diameters, bounded area ratios,
equal-area cells, WGS84 geodesics, or usable constant-time
coordinate-to-cell conversion. A proposed 64-bit layout
(version 4, resolution 4, root 6, path 50 bits) holds at most
ten base-32 refinement digits but is not an indexing library.
No timing benchmark or H3-vs-V5 speedup is claimed.

\section{Conclusion}
The narrow weak-existence question has a constructive affirmative
answer, under explicit nonconvexity and hanging-node allowances.
The geometric and algorithmic conditions required for an
operational pentagonal DGGS remain research questions.
\begin{thebibliography}{9}
\bibitem{h3docs} Uber Technologies, H3 Documentation,
\url{https://h3geo.org/docs/}.
\bibitem{h3code} Uber Technologies, H3 Source Code,
\url{https://github.com/uber/h3}.
\bibitem{ogc} Open Geospatial Consortium, Discrete Global Grid
Systems (DGGS), \url{https://docs.ogc.org/as/20-040r3/20-040r3.html}.
\bibitem{classic} Classical results: Euler's polyhedral relation,
triangulation of simple polygons, and Girard's spherical excess theorem.
\end{thebibliography}
\end{document}
